\section{Modern Code Review and Its Automation}
\label{sec:bg-code-review}

Modern code review is the lightweight, asynchronous inspection of a proposed
change before it enters a shared branch, organised around pull requests on
hosted platforms~\cite{davila2021mcrslr}. Bacchelli and Bird found that
developers enter review expecting to find defects but most reliably obtain
knowledge transfer and awareness of alternative solutions; the gap they
identified is \emph{understanding}, meaning a reviewer who does not know what
depends on a piece of code cannot tell whether changing it is
safe~\cite{bacchelli2013}. Subsequent work confirms the pattern: useful comments
identify functional defects and consequences the author missed, while comments
about naming and formatting are judged unhelpful~\cite{bosu2015}; at scale,
review is fast and habitual, calibrated to maintain norms as much as to catch
bugs~\cite{sadowski2018}. These studies, together with the ISO/IEC~25010
quality model~\cite{iso25010}, supply the categories from which the evaluation
rubric in Chapter~\ref{chap:edas} is derived.

Automating this task has been framed as transduction from a diff to a comment.
CodeReviewer pre-trains on review data and fine-tunes for comment
generation~\cite{li2022codereviewer}; Tufano et al.\ show pre-trained models
outperform earlier baselines~\cite{tufano2022codereview}; general-purpose
instruction-following models subsequently displaced task-specific
fine-tuning~\cite{hou2023llmse,fan2023llmse}.

The same shift enabled a wave of commercial AI review tools. CodeRabbit's
code-graph analysis resolves cross-file dependencies and symbol definitions
to enrich each review, alongside the output of more than forty integrated
linters and static analysers~\cite{coderabbit2026}. Greptile builds a graph of the
repository's files, functions, and dependencies and reviews each change
against it~\cite{greptile2026}. Qodo~Merge (built on the open-source
PR-Agent) deploys specialised review agents that gather context from the
whole codebase, including dependent repositories~\cite{qodo2026}. GitHub's Copilot code-review feature was used by
over one million developers within a month of its public preview and reached
general availability in April 2025, while operating primarily on the
diff~\cite{copilotcr2025}. These systems demonstrate that
practitioners value repository-wide context, but none exposes the reasoning
to controlled evaluation: they are closed-source products, not matched
experimental arms.

Two industrial studies begin to close that gap. Cihan et al.\ deploy Qodo's
PR-Agent at Beko on 1,568 of 4,335 pull requests and find that 73.8\% of
automated comments are resolved, although average closure time
increases~\cite{cihan2025automated}. Sun et al.\ report 75\% precision and
12\,000 weekly users for ByteDance's BitsAI-CR
pipeline~\cite{sun2025bitsaicr}.

In every case the input is the change, and everything the model knows about
the rest of the repository comes from pre-training or from the tool's own
retrieval pipeline. Recent benchmarks treat that as the limiting factor.
ContextCRBench tests review with enriched
context~\cite{hu2025contextcrbench}. SWR-Bench assembles a thousand verified
pull requests~\cite{zeng2025swrbench}. CodeFuse-CR-Bench evaluates at
repository level~\cite{guo2025codefusecrbench}. MCR-Bench benchmarks
multi-round review across 2,269 tasks and finds that 23\% of missed defects
are missed because the consequence lies in a different
file~\cite{zheng2026mcrbench}. The cross-file gap that MCR-Bench quantifies
is the blind spot targeted with a knowledge graph in this thesis.

While these benchmarks measure the effect of context, a parallel line of work
combines static-analysis output directly with LLM review. Jaoua et al.\
integrate PMD and Checkstyle diagnostics with a fine-tuned CodeLlama via
three strategies and find that the hybrid improves relevance and
completeness~\cite{jaoua2025combining}. In vulnerability detection, LLMxCPG
uses Joern to extract code-property-graph slices that reduce code volume by
67--91\% while preserving vulnerability-relevant context, improving F1 by
15--40\% over prior baselines~\cite{lekssays2025llmxcpg}. These results show that
static-analysis artefacts can usefully condition an LLM, but neither applies
CPG-level structural context to pull-request review. Jaoua et al.\ supply
linter diagnostics (rule violations), not a structural graph; LLMxCPG targets
security auditing, not review-comment generation.

All of these benchmarks and combination strategies establish that context
matters but do not compare structural and similarity-based selection within
one matched review pipeline, which is the comparison this thesis runs.
